\documentclass{beamer}
\usetheme{Madrid} % Ein professionelles Standard-Theme
\usecolortheme{whale}

\usepackage[utf8]{inputenc}
\usepackage[T1]{fontenc}
\usepackage[ngerman]{babel}
\usepackage{amsmath}
\usepackage{amsfonts}
\usepackage{amssymb}

\title{Determinismus und die Geometrie der Zahlen}
\subtitle{Gödels Grenzen vs. das Quaternionen-Primzahlmodell (\#Energiedoku)}
\author{Ihr Name / \#Energiedoku}
\date{\today}

\begin{document}
	
	% Titelseite
	\begin{frame}
		\titlepage
	\end{frame}
	
	% Inhaltsverzeichnis
	\begin{frame}{Überblick}
		\tableofcontents
	\end{frame}
	
	\section{Gödels Ende vs. Neue Strukturen}
	\begin{frame}{1. Gödels Ende vs. Die neue Struktur}
		Gödels Unvollständigkeitssatz markierte das Ende des naiven Determinismus.
		\begin{itemize}
			\item \textbf{Vom Chaos zur Ordnung:} Während Primzahlen in der klassischen Theorie oft ungeordnet wirken, bietet das Modell der \textbf{Quaternionen-Primzahlen} eine algebraische Struktur in $\mathbb{H}$.
			\item \textbf{Struktur statt Beweisbarkeit:} Gödels Limitierung der Axiomensysteme bedeutet nicht die Abwesenheit von Determinismus, sondern weist auf die Unvollständigkeit unserer logischen Werkzeuge hin.
			\item \textbf{Die Hypothese:} Die Welt ist deterministisch, doch unsere klassische Logik ist oft zu niederdimensional für die zugrunde liegende Realität.
		\end{itemize}
	\end{frame}
	
	\section{Die Welt als Rechenvorgang}
	\begin{frame}{2. Die Welt als algorithmischer Prozess}
		Die Verbindung von Computer Science und VWL führt zu einer Sichtweise der \textbf{algorithmischen Präzision}.
		\begin{block}{Quaternionen als Brücke}
			Durch die Nutzung von vier Dimensionen:
			\[ q = a + bi + cj + dk \]
			erlauben Quaternionen die Beschreibung von Rotationen und Energiezuständen, die in der reellen Ebene als stochastisches Rauschen erscheinen würden.
		\end{block}
		\begin{itemize}
			\item \textbf{Deterministisches Chaos:} Komplexe Systeme folgen präzisen, wenn auch hochsensiblen Regeln.
		\end{itemize}
	\end{frame}
	
	\section{Riemannsche Vermutung}
	\begin{frame}{3. Die Riemannsche Vermutung als Schlüssel}
		Die Verteilung der Primzahlen könnte das Ergebnis eines \textbf{schwingenden, deterministischen Systems} sein.
		\begin{itemize}
			\item \textbf{Polya-Hilbert-Vermutung:} Nullstellen der Zeta-Funktion als Eigenwerte physikalischer Operatoren.
			\item \textbf{Geometrische Architektur:} Wir bewegen uns weg von Wahrscheinlichkeiten hin zur Architektur der Zahlen.
		\end{itemize}
		\begin{quote}
			Gödels „Dunkelheit“ ist der Schatten einer vierdimensionalen Wahrheit (Quaternionen) in einem eindimensionalen Logikgitter.
		\end{quote}
	\end{frame}
	
	\section{Fazit}
	\begin{frame}{Fazit der \#Energiedoku}
		Liefert die Verbindung von Hamiltons Quaternionen mit der analytischen Zahlentheorie das nötige Licht, um Gödels Schatten zu durchbrechen?
		
		\vspace{1cm}
		\begin{center}
			\textbf{Frage zur Diskussion:} \\
			Kann die physikalische Realität der Primzahlen durch die Erweiterung des Zahlenraums vollständig determiniert werden?
		\end{center}
	\end{frame}
	
\end{document}